\newpage
\section*{Résumé}

L'asservissement visuel classique nécessite une extraction et un suivi robuste des caractéristiques géométriques (points, lignes, contours) pour construire la matrice d'interaction reliant les mouvements de la caméra aux variations observées dans l'image. Cette approche est efficace mais fortement  dépendante de la qualité de l'extraction de ces caractéristiques et difficile à généraliser pour des environnements complexes  ou réels. L'asservissement visuel photométrique s'affranchit de cette étape en exploitant directement l'intensité des pixels comme signal de commande, mais devient alors sensible à la dimensionnalité de l'espace image et aux minima locaux de la fonction de coût.
Le  travail effectué pendant ce stage propose de combiner l'asservissement visuel photométrique, formulé dans un espace latent de faible dimension, avec l'apprentissage par renforcement (RL) pour la navigation visuelle. L'objectif est d'exploiter les capacités de généralisation et de prise de décision séquentielle du RL pour piloter un agent dans un espace de représentation latente construit à partir de l'information photométrique. Les environnements de navigation utilisés pour entraîner et évaluer les modèles  sont reconstruits par Gaussian Splatting (3DGS), une méthode de reconstruction 3D permettant de générer, à partir d'un ensemble d'images, des vues photoréalistes denses et continues d'une scène, offrant ainsi un cadre photométrique riche et cohérent pour l'apprentissage de l'asservissement visuel.

\vspace{0.5cm}
\textbf{Mots-clés :} asservissement visuel photométrique, espace latent, apprentissage par renforcement profond, Gaussian Splatting, navigation visuelle.

